
\documentclass{article}
\usepackage{colortbl}
\usepackage{makecell}
\usepackage{multirow}
\usepackage{supertabular}

\begin{document}

\newcounter{utterance}

\centering \large Interaction Transcript for game `eqbench', experiment `eq\_bench', episode 57 with reflexion{-}Qwen3.5{-}27B.
\vspace{24pt}

{ \footnotesize  \setcounter{utterance}{1}
\setlength{\tabcolsep}{0pt}
\begin{supertabular}{c@{$\;$}|p{.15\linewidth}@{}p{.15\linewidth}p{.15\linewidth}p{.15\linewidth}p{.15\linewidth}p{.15\linewidth}}
    \# & \multicolumn{2}{c}{Player} && \multicolumn{2}{c}{Game Master} \\
    \hline

    \theutterance \stepcounter{utterance}  
    & & & \multicolumn{4}{p{0.6\linewidth}}{
        \cellcolor[rgb]{0.9,0.9,0.9}{
            \makecell[{{p{\linewidth}}}]{
                \texttt{\tiny{[P1$\langle$GM]}}
                \texttt{You have attempted this same task before and did not succeed. These are the plans you wrote after those attempts.} \\
\texttt{Trial \#1: New plan: My previous attempt failed because I assigned a high score to "Defensiveness" (9) while overlooking the stronger textual evidence for "Regret." Tyler's accusation ("You always do, don't you?") and Claudia's incomplete, hesitant response ("I had to...") suggest she is grappling with the consequences of her past actions rather than actively defending herself. The phrase "I try to understand" indicates an attempt at empathy, which aligns more with regret over the hurt caused than with defensiveness. Next time, I will prioritize the emotion that best explains the character's hesitation and the specific context of "leaving someone alone," which points to guilt and regret. I will lower the score for Defensiveness and significantly increase the score for Regret to reflect the emotional weight of the unresolved departure.} \\
\\ 
\texttt{Trial \#2: New plan: My previous attempt failed because I underestimated the intensity of "Defensiveness" relative to "Regret." While Claudia's hesitation ("I had to...") suggests regret, her immediate reaction to Tyler's accusation ("That's not fair, Ty") is a classic defensive maneuver to deflect blame. Tyler's aggressive questioning ("Chase your dreams? Leave your past behind?") forces her into a corner where she must justify her actions, making defensiveness the most immediate and intense reaction. Next time, I will prioritize "Defensiveness" as the highest score (e.g., 9) because her dialogue explicitly rejects the fairness of the accusation. I will lower "Regret" to a moderate score (e.g., 5) to acknowledge the underlying guilt without letting it overshadow the active defense. "Confusion" and "Pride" will remain low, as the dialogue shows clear understanding of the conflict and no indication of pride.} \\
\\ 
\\ 
\\ 
\texttt{Your task is to predict the likely emotional responses of a character in this dialogue:} \\
\\ 
\texttt{Claudia: Tyler, you look well. Teaching suits you.} \\
\texttt{Tyler: Yeah, well, someone's gotta do the dirty work.} \\
\texttt{Claudia: I didn't mean it like that. I...} \\
\texttt{Tyler: I know what you meant. You always do, don't you?} \\
\texttt{Claudia: Not always, no. But I try to understand, at least.} \\
\texttt{Tyler: Understand? Like you understood when you left me alone?} \\
\texttt{Claudia: That's not fair, Ty. I had to...} \\
\texttt{Tyler: What? Chase your dreams? Leave your past behind?} \\
\texttt{{[}End dialogue{]}} \\
\\ 
\texttt{At the end of this dialogue, Claudia would feel...} \\
\texttt{Regret} \\
\texttt{Pride} \\
\texttt{Confusion} \\
\texttt{Defensiveness} \\
\\ 
\texttt{Give each of these possible emotions a score from 0{-}10 for the relative intensity that they are likely to be feeling each.} \\
\\ 
\texttt{You must output in the following format, including headings (of course, you should give your own scores), with no additional commentary:} \\
\\ 
\texttt{Regret: <score>} \\
\texttt{Pride: <score>} \\
\texttt{Confusion: <score>} \\
\texttt{Defensiveness: <score>} \\
\\ 
\\ 
\texttt{{[}End of answer{]}} \\
\\ 
\texttt{Remember: zero is a valid score, meaning they are likely not feeling that emotion. You must score at least one emotion > 0.} \\
\\ 
\texttt{Your answer:} \\
            }
        }
    }
     \\ \\

    \theutterance \stepcounter{utterance}  
    & \multicolumn{4}{p{0.6\linewidth}}{
        \cellcolor[rgb]{0.561,0.737,0.561}{
            \makecell[{{p{\linewidth}}}]{
                \texttt{\tiny{[P1$\rangle$GM]}}
                \texttt{Regret: 8} \\
\texttt{Pride: 1} \\
\texttt{Confusion: 2} \\
\texttt{Defensiveness: 9} \\
            }
        }
    }
    & & \\ \\

    \theutterance \stepcounter{utterance}  
    & & & \multicolumn{2}{p{0.3\linewidth}}{
        \cellcolor[rgb]{0.9,0.9,0.9}{
            \makecell[{{p{\linewidth}}}]{
                \texttt{\tiny{[GM$|$GM]}}
                \texttt{Regret: 8} \\
\texttt{Pride: 1} \\
\texttt{Confusion: 2} \\
\texttt{Defensiveness: 9} \\
            }
        }
    }
    & & \\ \\

    \theutterance \stepcounter{utterance}  
    & & & \multicolumn{2}{p{0.3\linewidth}}{
        \cellcolor[rgb]{0.9,0.9,0.9}{
            \makecell[{{p{\linewidth}}}]{
                \texttt{\tiny{[GM$|$GM]}}
                \texttt{Regret: 8} \\
\texttt{Pride: 0} \\
\texttt{Confusion: 0} \\
\texttt{Defensiveness: 6} \\
            }
        }
    }
    & & \\ \\

    \theutterance \stepcounter{utterance}  
    & & & \multicolumn{2}{p{0.3\linewidth}}{
        \cellcolor[rgb]{0.9,0.9,0.9}{
            \makecell[{{p{\linewidth}}}]{
                \texttt{\tiny{[GM$|$GM]}}
                \texttt{game\_result = LOSE} \\
            }
        }
    }
    & & \\ \\

    \theutterance \stepcounter{utterance}  
    & & \multicolumn{4}{p{0.6\linewidth}}{
        \cellcolor[rgb]{0.95,0.95,0.95}{
            \makecell[{{p{\linewidth}}}]{
                \texttt{\tiny{[P1$\leftrightarrow$P]}}
                \texttt{Post{-}episode Reflexion (not sent to GM):} \\
\texttt{New plan: My previous attempts failed because I treated "Regret" and "Defensiveness" as mutually exclusive or prioritized one based on a single line of dialogue, leading to inconsistent scoring across trials. I need to recognize that Claudia is likely experiencing both emotions simultaneously but with different intensities based on the dialogue's progression. Her initial "That's not fair" is defensive, but her trailing off ("I had to...") and the context of Tyler's hurt ("left me alone") reveal deep underlying regret. Instead of swinging between high scores for one or the other, I will assign high scores to *both* to reflect the complexity of her state: high Defensiveness (8) for her immediate reaction to the accusation, and high Regret (8) for the emotional weight of her past actions. I will ensure "Pride" remains near zero (0) as there is no evidence of it, and "Confusion" remains low (1) since she clearly understands the situation but is struggling to articulate her justification. This balanced approach accounts for the full emotional arc of the scene.} \\
            }
        }
    }
    & \\ \\

\end{supertabular}
}

\end{document}
